\section{Introduction}
\label{sec:intro}

% Introduce 4D gaussian splatting
% Closed-loop evaluation has become increasingly important for end-to-end autonomous driving.
% Such an evaluation requires the ego vehicle to receive new observations that correspond to its actions in the simulated environment.
% Visual scene reconstruction and novel view synthesis from real urban scenes, which are captured by autonomous vehicle sensors, particularly multi-view cameras, can serve as the foundation for achieving realistic closed-loop evaluation in end-to-end autonomous driving systems.
% There are three common ways to achieve this goal: extracting and importing real elements into simulators such as CARLA~\cite{dosovitskiy2017carla} and rendering them using game engines; implicitly representing and generating scenes through models such as NeRF~\cite{mildenhall2021nerf} or video diffusion models; and explicitly representing scenes with 3D geometric structures.
% Among these methods, 4D Gaussian Splatting (4DGS)~\cite{wu20244d}, an extension of 3D Gaussian Splatting (3DGS)~\cite{kerbl20233d} with an additional temporal dimension, effectively balances geometric accuracy, photometric fidelity, and real-time rendering efficiency. This makes it well-suited for dynamic and interactive simulated evaluation in end-to-end autonomous driving.

Closed-loop evaluation is increasingly vital for end-to-end autonomous driving, requiring the ego-vehicle to receive sensor observations that faithfully correspond to its actions within a simulated environment. Visual scene reconstruction and novel-view synthesis from real-world multi-view imagery provide the bedrock for such simulations. Existing approaches generally fall into three categories: asset importation into game engines like CARLA~\cite{dosovitskiy2017carla}, implicit representation via NeRF~\cite{mildenhall2021nerf} or video diffusion, and explicit geometric modeling. Among these, 4D Gaussian Splatting (4DGS)~\cite{wu20244d}, an extension of 3DGS~\cite{kerbl20233d} with a temporal dimension, effectively balances geometric accuracy, photometric fidelity, and real-time rendering, making it ideal for interactive driving simulations.

% Challenges
% Traditional 4D Gaussian Splatting (4DGS) relies on per-scene optimization, as seen in methods like Street Gaussian~\cite{yan2024street}. These approaches reconstruct scenes by iteratively optimizing Gaussian kernels to minimize rendering loss against ground-truth images, often requiring geometry priors like LiDAR point clouds for initialization. However, this non-data-driven paradigm fails to leverage shared knowledge across scenes and demands substantial computational overhead for every new environment, limiting its scalability.
% In contrast, the success of feed-forward models like VGGT~\cite{wang2025vggt}, which directly produce explicit 3D geometry such as per-pixel depth and point maps, demonstrates the potential of 3D foundation models for efficient reconstruction.
% However, directly extending VGGT~\cite{wang2025vggt} to 4D Gaussian reconstruction for autonomous driving presents several challenges:
% 1) Photometric Deficiency. While VGGT recovers 3D geometry, its features lack the fine-grained photometric detail (e.g., color and texture) necessary to regress Gaussian appearance attributes for high-fidelity novel-view synthesis.
% 2) Temporal Staticity. As a static 3D foundation model, VGGT lacks mechanisms to model the temporal dimension, making it incapable of representing the dynamic motion of traffic participants inherent in 4D driving scenes.
% 3) Domain and Calibration Mismatch. Significant spatial errors occur when applying VGGT's pre-trained weights to urban driving scenes due to the data domain gap. Furthermore, VGGT treats camera parameters as unknowns, failing to utilize the pre-calibrated information in driving datasets.
However, traditional 4DGS methods, such as Street Gaussian~\cite{yan2024street}, rely on per-scene optimization. These approaches iteratively optimize Gaussian kernels to minimize rendering loss, often requiring LiDAR priors for initialization. This non-data-driven paradigm fails to leverage shared structural knowledge across scenes and incurs substantial computational overhead for every new environment, limiting scalability. In contrast, feed-forward models like VGGT~\cite{wang2025vggt} demonstrate the potential of 3D foundation models for efficient reconstruction by directly producing explicit geometry. Yet, extending VGGT to 4D Gaussian reconstruction for autonomous driving presents three primary challenges. First is photometric deficiency, as VGGT features lack the fine-grained detail necessary to regress high-fidelity appearance attributes. Second is temporal staticity, where static backbones cannot represent the dynamic motion of traffic participants. Third is domain and calibration mismatch, where data gaps lead to 3D geometry prediction errors and a failure to utilize the pre-calibrated sensor intrinsics and extrinsics inherent in driving datasets.


\begin{figure*}[t]
    \setlength{\belowcaptionskip}{0pt}
	\centering
	\includegraphics[width=1.0\textwidth]{figures/recondrive-framework.pdf}
  \vspace{-9pt}
	\caption{\textbf{ReconDrive Inference Framework.} ReconDrive is a powerful feed-forward 4D Gaussian splatting generation framework tailored for urban scene reconstruction and novel-view synthesis. In this framework, we select two context frames from each segment of an urban scene as input, and adopt a static-dynamic composition strategy to represent 4D Gaussians. To further enhance geometric precision, we design a dedicated Gaussian prediction head, consisting of a Gaussian Parameter Prediction Head (GPPH) and a Gaussian Center Prediction Head (GCPH), which enables the generation of Gaussians with more accurate geometric attributes.}
 \label{fig:ReconDrive-framework}
     \vspace{-8pt}
\end{figure*}

To address these limitations, we propose ReconDrive, a feed-forward framework that reconstructs 4D Gaussian Splatting representations from vision-only inputs. ReconDrive utilizes a pre-trained VGGT-style backbone to directly regress spatio-temporal Gaussian parameters, enabling efficient reconstruction at scale without per-scene optimization. The inference pipeline is built upon three design pillars. First, we employ hybrid Gaussian prediction heads; we concatenate raw images with dense pixel-wise features to provide the photometric cues required for regressing appearance attributes like opacity and spherical harmonics, while incorporating camera calibration directly into the centers head for precise spatial positioning. Second, we utilize a static-dynamic 4D composition, leveraging the SAM2~\cite{ravi2024sam} foundation model to extract object masks and assign temporal velocities to dynamic Gaussians through pixel-wise correspondence. Third, we implement segment-wise temporal fusion to handle view sparsity over long driving sequences by partitioning the scene into segments and fusing localized Gaussian clusters into a unified 4D representation.
The overall inference framework is illustrated in Fig.~\ref{fig:ReconDrive-framework}.
During training, we utilize frozen VGGT~\cite{wang2025vggt} weights and apply LoRA~\cite{hu2022lora} for parameter-efficient fine-tuning on urban driving data. We introduce a novel-frame projection loss to facilitate stable geometry reconstruction against unseen temporal views.

To ensure a fair comparison, we establish a comprehensive benchmark on the nuScenes~\cite{caesar2020nuscenes} dataset, reproducing five representative optimization and feed-forward methods under a unified protocol. We evaluate performance across scene reconstruction, novel-view synthesis, and downstream tasks such as 3D object detection and tracking. Within this benchmark, ReconDrive significantly outperforms existing feed-forward approaches across all metrics and surpasses per-scene optimization methods in eight out of nine categories. Our results demonstrate that feed-forward approaches can achieve competitive and often superior novel-view synthesis performance while requiring substantially fewer resources for Gaussian generation. This highlights the promising potential of feed-forward paradigms for large-scale autonomous driving reconstruction.

% Contributions
Our contributions can be summarized as follows.
\begin{itemize}[leftmargin=9pt]
\item We introduce ReconDrive, a feed-forward framework that can directly generate 4D Gaussian splatting without per-scene optimization, to reconstruct urban scenes and enable novel-view synthesis across time.
\item We introduce hybrid Gaussian prediction heads, a static-dynamic composition strategy, and a segment-wise temporal fusion representation. Combined with specialized training, these components enable 3D foundation models in dynamic urban environments.
\item We build autonomous driving scene reconstruction benchmark by reproducing typical baselines on the nuScenes~\cite{caesar2020nuscenes}. Our ReconDrive achieves state-of-the-art performance among feed-forward methods and competitive results against optimization-based baselines across multiple evaluation protocols.
\end{itemize}